% ============================================================
% 第 20 章 直流电机概述（教材第六篇，印刷页 227--234）
% 节名：第一节 直流电机的工作原理；第二节 直流电机的结构；第三节 直流电机的额定值
% ============================================================
\chapter{直流电机概述}

直流电机是历史上最早实用的电机，至今仍在需要宽范围平滑调速、大起动转矩的场合（轧钢、矿山、电车、起重）中大量使用。它的核心是一个\textbf{机械换向器}：把电枢绕组内部的交流电动势"整流"成对外输出的直流，或把外部直流"逆变"成绕组内部的交流。本章介绍原理、结构与额定值。

\section{直流电机的工作原理}

\subsection{直流电动机}
电枢绕组中的导体位于定子主磁极产生的磁场中，通入电流后受电磁力 $f=Bli$ 作用产生转矩，使电枢旋转。若导体固定不动，转矩方向恒定；但导体随转子旋转到另一极下时，磁场方向变了，若不改变电流方向，转矩就会反向，电机无法连续旋转。

解决的办法是\textbf{换向器}：它由许多彼此绝缘、与电枢绕组线圈按一定规律连接的换向片组成，电刷固定不动压在换向器表面。线圈从一个极下转到另一个极下时，电刷使线圈中的电流方向同时改变，因此\textbf{转矩方向始终保持不变}。

因此，换向器的作用是：\textbf{对电动机，把外部直流"逆变"为电枢绕组内部的交流（不断改向的电流）；对发电机，把绕组内部的交流电动势"整流"为对外输出的直流。}换向器称为\textbf{机械式逆变器/整流器}。

\subsection{直流发电机}
原动机拖动电枢旋转，导体切割主磁通感应电动势 $e=Blv$，方向由右手定则确定；同样由于导体交替进入不同极性的磁极下，绕组内部电动势是交变的，经换向器与电刷才输出方向不变的直流电动势。

\subsection{可逆性}
同一台直流电机，通入直流电流产生转矩、拖动机械负载，即为电动机；由原动机拖动产生电动势、输出直流电能，即为发电机。故直流电机同样具有可逆性，与第 2 章"能量守恒与机电能量转换"的结论一致。

\subsection{换向器工作的分步描述}
以最简单的单匝线圈为例（线圈两端分别接两个换向片，电刷压在换向片上）：
\begin{enumerate}
  \item 线圈边 $ab$ 在 $N$ 极下、$cd$ 在 $S$ 极下，电流由 $a\to b\to c\to d$，两导体所受电磁力方向一致，产生同一方向的转矩；
  \item 电枢转过 $180^\circ$ 后，$ab$ 到了 $S$ 极下、$cd$ 到了 $N$ 极下；若电流方向不变，转矩就会反向；
  \item 恰在此刻，两换向片被同一电刷短接，线圈中电流换向；随后电刷接触的换向片互换，线圈中电流反向；
  \item 于是 $ab$ 在 $S$ 极下时电流方向已反过来，转矩方向仍与原来一致——电枢得以连续旋转。
\end{enumerate}
这个过程每转半周重复一次，故电枢绕组内部的电流是\textbf{交变}的，而对外（电刷端）的电流方向恒定。这就是"机械换向"的本质：\textbf{靠机械接触的位置变化，在电流过零的瞬间完成换向}。

\subsection{与交流电机的对比}
\begin{center}
\begin{tabular}{@{}p{2.4cm}p{4.4cm}p{4.4cm}@{}}
\toprule
项目 & 直流电机 & 交流电机 \\
\midrule
磁场 & 定子主磁极（直流励磁） & 定子绕组产生旋转磁场 \\
电枢 & 转子（旋转） & 定子（静止） \\
绕组内电流/电动势 & 交流（靠换向器整流/逆变） & 交流（本身就是交流） \\
对外电流 & 直流 & 交流 \\
薄弱环节 & 换向器与电刷 & 无（笼型） \\
\bottomrule
\end{tabular}
\end{center}

\section{直流电机的结构}

直流电机由定子（静止部分）和转子（电枢，旋转部分）组成，两者之间有气隙。

\subsection{定子部分}
\begin{itemize}
  \item \textbf{主磁极}：由磁极铁芯（薄钢板叠成）和套在其上的\textbf{励磁绕组}构成，用于建立主磁场。磁极下部的扩大部分称\textbf{极靴}，使气隙磁密分布合理。
  \item \textbf{换向极}：装在相邻主磁极之间（几何中性线上），绕组与电枢绕组串联，用于改善换向（见第 25 章）。
  \item \textbf{机座}：既作机械支撑，又作为主磁路的一部分（磁轭）。
  \item \textbf{电刷装置}：由电刷、刷握、刷杆和连线组成，把电枢绕组与外电路连接起来。
\end{itemize}

\subsection{转子（电枢）部分}
\begin{itemize}
  \item \textbf{电枢铁芯}：由硅钢片叠成，外圆冲有齿槽以嵌放电枢绕组，同时构成主磁路。
  \item \textbf{电枢绕组}：由绝缘铜线绕成线圈嵌入槽内，按一定规律与换向片连接（单叠或单波，见第 21 章）。
  \item \textbf{换向器}：由许多梯形铜换向片与片间云母绝缘组成，是直流电机的关键部件，也是最易出故障的部件。
  \item \textbf{转轴、风扇}：支撑与冷却。
\end{itemize}

\begin{center}
\begin{tabular}{@{}p{2.2cm}p{4.6cm}p{5.2cm}@{}}
\toprule
部件 & 作用 & 备注 \\
\midrule
主磁极 & 建立主磁场 & 励磁绕组通直流 \\
换向极 & 改善换向、抑制火花 & 绕组与电枢串联 \\
机座 & 机械支撑 + 磁轭 & 构成磁路的一部分 \\
电枢铁芯 & 嵌放绕组 + 主磁路 & 硅钢片叠压 \\
电枢绕组 & 感应电动势、产生转矩 & 经换向片连接 \\
换向器 & 机械整流/逆变 & 直流电机的核心与薄弱环节 \\
电刷 & 电枢与外电路的滑动连接 & 需定期更换 \\
\bottomrule
\end{tabular}
\end{center}

\section{直流电机的额定值}

直流电机的额定值都标注在铭牌上，是电机长期安全运行的规定值。

\subsection{额定功率 $P_N$}
\textbf{电动机}：额定运行时轴上输出的机械功率（kW）；\textbf{发电机}：额定运行时出线端输出的电功率（kW）。注意两者定义不同。

\subsection{额定电压 $U_N$}
额定运行时的端电压（V）。直流电机额定电压等级常见为 110、220、440、750\,V 等。

\subsection{额定电流 $I_N$}
额定运行时出线端的电流（A）。对电动机，$I_N$ 是输入电流；对发电机，$I_N$ 是输出电流。

\subsection{额定转速 $n_N$}
额定运行时的转速（r/min）。

\subsection{额定励磁电压 $U_{fN}$ 与额定励磁电流 $I_{fN}$}
额定运行时的励磁电压与励磁电流。

\subsection{额定效率 $\eta_N$}
额定运行时输出功率与输入功率之比。

\begin{center}
\begin{tabular}{@{}lll@{}}
\toprule
额定值 & 电动机 & 发电机 \\
\midrule
$P_N$ & 轴上输出机械功率 & 出线端输出电功率 \\
$U_N$ & 输入端电压 & 输出端电压 \\
$I_N$ & 输入电流 & 输出电流 \\
基本关系 & $P_N=U_NI_N\eta_N$ & $P_N=U_NI_N$ \\
\bottomrule
\end{tabular}
\end{center}

\keyformula{\text{电动机：}P_1=U_NI_N,\ P_N=U_NI_N\eta_N;\qquad
\text{发电机：}P_2=U_NI_N,\ P_1=\frac{U_NI_N}{\eta_N}}

\corepoint{直流电机额定值最容易混的是 $P_N$ 的定义：电动机的 $P_N$ 是\textbf{输出机械功率}，发电机的 $P_N$ 是\textbf{输出电功率}。因此同样铭牌"$P_N=10$\,kW、$U_N=220$\,V"，电动机的额定电流是 $10\,000/(220\eta_N)$，发电机的额定电流是 $10\,000/220$——一个要除效率、一个不除。做题先问"这是电动机还是发电机"。}
